\section{Health and network}
\label{sec:network}

\begin{question}
  \textbf{Are all GPUs and links alike, and does the fabric have a structure?} Do all 416 B200 deliver the
  same matmul rate? Does every node reach its NVLink and InfiniBand bandwidth? Do some node pairs
  communicate faster than others? Does collective bandwidth hold up to 52 nodes?
\end{question}

This group runs first: it excludes faulty nodes, and it decides whether node placement is a factor of
the later experiments.

\paragraph{GPUs (D1).} A per-node health check measures, on every GPU, the BF16 and FP8 matmul rate, and
on every node the NVLink all-reduce bandwidth, the loopback bandwidth of each HCA and the GPU--HCA
affinity. \Cref{fig:health_gpus} compares every GPU with the median.

\begin{figure}[h]
  \centering\fig{health_gpus}
  \caption{Matmul rate of every GPU in BF16 and FP8, by node (strip), with the median and a line 5\,\%
    below it. \syntheticnote{health_gpus}}
  \label{fig:health_gpus}
\end{figure}

\placeholder{median rates against SciNet's numbers; GPUs below the $-5\,\%$ line and what was done with them.}

\paragraph{Links (D1).} \Cref{fig:health_links} shows the intra-node and per-HCA bandwidths.

\begin{figure}[h]
  \centering\fig{health_links}
  \caption{Per-node NVLink all-reduce bus bandwidth (left) and per-HCA bandwidth (right).
    \syntheticnote{health_links}}
  \label{fig:health_links}
\end{figure}

\placeholder{number of HCAs per node and their rate; outlier nodes or links.}

\paragraph{Pairs (D2).} Every pair of nodes is measured: 26 disjoint pairs at a time, 51 rounds, all
1{,}326 pairs. If the matrix shows groups of nodes (latency or bandwidth clusters, for example along the
8 subnets of \texttt{/etc/hosts}), a placement experiment is added (Qwen and GLM-5 at 13 nodes inside one
group vs.\ spread over groups). Otherwise placement is dropped from the design.

\begin{figure}[h]
  \centering\fig{pair_matrix}
  \caption{Pairwise bandwidth (and latency) between the 52 nodes, ordered by subnet.
    \syntheticnote{pair_matrix}}
  \label{fig:pair_matrix}
\end{figure}

\placeholder{whether the matrix has a structure, and the decision on placement.}

\paragraph{Collectives at scale (D3, D4).} NCCL all-reduce, all-gather, reduce-scatter and all-to-all
with 1\,GiB messages at 1, 2, 4, 8, 13, 26 and 52 nodes, and a message-size sweep at 2 and 52 nodes. The
bisection test (D4) splits the nodes into two random halves of 26 that exchange data all-to-all, for
3 random cuts. A non-blocking fabric keeps the per-node bandwidth constant as nodes are added.

\begin{figure}[h]
  \centering\fig{collectives_scale}
  \caption{Bus bandwidth vs.\ number of nodes, one line per collective, with the 800\,Gb/s line rate of a
    node; bisection points (D4). \syntheticnote{collectives_scale}}
  \label{fig:collectives_scale}
\end{figure}

\placeholder{bandwidth at 52 nodes as a share of the line rate; bisection bandwidth; whether the fabric
behaves as non-blocking.}
